\chapter{Esophagectomy}

\section*{What Is This Surgery?}
Esophagectomy is a complex and demanding surgical procedure involving the partial or total removal of the esophagus, the muscular tube connecting the pharynx to the stomach. This major operation is primarily performed for the curative treatment of esophageal cancer, but it is also indicated for certain end-stage benign conditions that severely impair esophageal function or pose a significant risk to the patient. Following resection, gastrointestinal continuity is typically restored by mobilizing the stomach or a segment of the colon or jejunum to serve as a conduit, which is then anastomosed to the remaining esophagus or pharynx. Given the esophagus's deep anatomical location within the mediastinum, adjacent to vital structures such as the heart, lungs, aorta, and trachea, esophagectomy carries substantial physiological stress and requires meticulous surgical technique, comprehensive preoperative optimization, and intensive postoperative care. The goal is to achieve oncologic clearance while preserving or restoring the patient's ability to swallow and maintain nutrition.

\section*{Alternative Names}
\begin{itemize}
  \item Total esophagectomy
  \item Partial esophagectomy
  \item Esophagogastrectomy
  \item Ivor-Lewis esophagectomy (right thoracotomy and laparotomy)
  \item Transhiatal esophagectomy (abdominal and cervical incisions)
  \item McKeown esophagectomy (three-incision or three-field: abdominal, right thoracotomy, cervical)
  \item Minimally invasive esophagectomy (MIE)
  \item Hybrid esophagectomy
\end{itemize}

\section*{Relevant Surgical Anatomy}
The esophagus is a muscular tube, approximately 25-30 cm long in adults, extending from the cricoid cartilage (C6 vertebral level) to the cardia of the stomach (T11 vertebral level). It traverses the neck, superior and posterior mediastinum, and a short segment within the abdomen. Its wall consists of four layers: mucosa, submucosa, muscularis propria (inner circular, outer longitudinal), and adventitia (no serosa in most of its length). Key anatomical relations include the trachea anteriorly in the neck and superior mediastinum, the aortic arch and left main bronchus anteriorly in the mid-thorax, and the pericardium and left atrium more inferiorly. Posteriorly, it lies against the vertebral column and thoracic duct.

The arterial supply to the esophagus is segmental and variable: the cervical esophagus receives branches from the inferior thyroid arteries; the thoracic esophagus from bronchial arteries, direct aortic branches, and intercostal arteries; and the abdominal esophagus from the left gastric artery and inferior phrenic arteries. Venous drainage largely parallels the arterial supply, with the cervical esophagus draining into the inferior thyroid veins, the thoracic esophagus into the azygos and hemiazygos veins, and the abdominal esophagus into the left gastric vein, which is part of the portal system. This portosystemic connection is clinically significant in portal hypertension. Lymphatic drainage is extensive and longitudinal, allowing for widespread tumor dissemination; it includes cervical, paratracheal, subcarinal, paraesophageal, celiac, and perigastric lymph nodes. The vagus nerves (right and left) descend with the esophagus, forming the esophageal plexus and then anterior and posterior vagal trunks, which are typically resected during esophagectomy, leading to gastric denervation. The recurrent laryngeal nerves ascend in the tracheoesophageal groove and are at risk of injury during cervical dissection.

\section{Surgical Principle \& Rationale}
The fundamental principle of esophagectomy for malignancy is the complete en bloc resection of the tumor-bearing esophagus, along with its associated regional lymph nodes, to achieve a margin-negative (R0) resection. This aims to eradicate both the primary tumor and potential microscopic metastatic disease within the lymphatic basin, thereby maximizing the chance of cure. The rationale for extensive lymphadenectomy stems from the high propensity of esophageal cancers to metastasize early to regional lymph nodes, often distant from the primary tumor. Following resection, the critical second principle is the restoration of gastrointestinal continuity, most commonly achieved by mobilizing the stomach into the chest or neck to create a gastric conduit. This conduit is then anastomosed to the remaining proximal esophagus or pharynx, allowing the patient to resume oral intake. Technical goals include achieving wide tumor-free margins, performing a thorough and anatomically appropriate lymphadenectomy, ensuring a well-vascularized and tension-free anastomosis to minimize the risk of anastomotic leak, and preserving sufficient conduit length and perfusion to reach the cervical or intrathoracic esophageal remnant.

\section{History \& Surgical Evolution}
Early attempts at esophagectomy were fraught with prohibitive mortality due to infection and anastomotic complications. The first successful esophagectomy with reconstruction was performed by Franz Torek in 1913, who resected a cervical esophageal carcinoma and created an external skin tube to connect the cervical esophagus to the stomach, a staged and highly morbid procedure. Significant progress was made in 1933 when Tetsuzo Ohsawa performed the first successful one-stage intrathoracic esophagectomy with esophagogastric anastomosis for esophageal cancer. The modern era of esophagectomy was greatly influenced by Ivor Lewis in the 1940s, who popularized the two-stage approach involving a laparotomy for gastric mobilization followed by a right thoracotomy for esophageal resection and intrathoracic anastomosis. In the 1970s, Mark Orringer introduced the transhiatal esophagectomy, which avoided thoracotomy by dissecting the esophagus bluntly from the abdomen and neck, performing a cervical anastomosis. The McKeown (three-field) approach, combining abdominal, thoracic, and cervical incisions with a cervical anastomosis, gained favor for its comprehensive lymphadenectomy and safer cervical anastomosis. More recently, the advent of minimally invasive techniques, including thoracoscopy and laparoscopy, and robotic-assisted surgery, has aimed to reduce the morbidity of this extensive operation while maintaining equivalent oncologic outcomes, marking a significant evolution in patient recovery and surgical precision.

\section{Clinical Indications}
\begin{itemize}
  \item Localized or locally advanced esophageal cancer (adenocarcinoma or squamous cell carcinoma) suitable for curative resection, often following neoadjuvant chemoradiation therapy.
  \item High-grade dysplasia (HGD) in Barrett's esophagus that is not amenable to endoscopic mucosal resection or ablation, indicating a high risk of progression to invasive cancer.
  \item Persistent or recurrent esophageal cancer after prior definitive non-surgical therapies, where surgical salvage is deemed feasible and offers a chance for cure.
  \item Extensive esophageal perforation or rupture that cannot be managed conservatively or endoscopically, leading to severe mediastinitis or empyema.
  \item End-stage benign esophageal diseases, such as intractable caustic strictures, severe achalasia refractory to multiple endoscopic and surgical treatments (e.g., Heller myotomy), or diffuse esophageal spasm causing profound dysphagia and malnutrition.
  \item Selected cases of esophageal leiomyoma or other benign tumors that are too large for endoscopic removal or cause significant symptoms.
  \item Rarely, for severe, intractable gastroesophageal reflux disease (GERD) with extensive esophageal damage or failed fundoplication, though this is a last resort.
\end{itemize}

\section*{Clinical Positioning}
For esophageal cancer, esophagectomy is a primary curative treatment when the disease is localized or locally advanced and deemed resectable after thorough staging. It is often integrated into a multimodal treatment strategy, frequently preceded by neoadjuvant chemotherapy, radiation therapy, or chemoradiation, which aims to downstage the tumor, treat micrometastases, and improve R0 resection rates and overall survival. In this context, surgery is a critical component of the primary curative intent, rather than a secondary or salvage measure. For benign esophageal diseases, esophagectomy is almost invariably a secondary or tertiary option, reserved for cases where less invasive endoscopic or medical treatments have failed, or when the disease has progressed to an end-stage condition causing severe symptoms and compromising quality of life. The decision to proceed with esophagectomy is made after comprehensive multidisciplinary evaluation, considering the patient's overall health, tumor characteristics, and the potential for significant morbidity and mortality associated with the procedure.

\section{Surgical Technique \& Procedure}
Esophagectomy is performed under general anesthesia, typically with a double-lumen endotracheal tube for single-lung ventilation during thoracic phases, and invasive monitoring including arterial and central venous access. Patient positioning varies by approach: supine for transhiatal, left lateral decubitus for right thoracotomy (Ivor-Lewis, McKeown thoracic phase), and often repositioning for combined approaches.

The general steps for a typical McKeown (three-field) esophagectomy, which includes abdominal, thoracic, and cervical phases, are as follows:

\begin{enumerate}
  \item  **Abdominal Phase (Laparotomy or Laparoscopy):**
    *   Upper midline incision or laparoscopic ports are placed.
    *   The stomach is fully mobilized, preserving the right gastroepiploic artery and vein as the primary blood supply for the gastric conduit. The left gastric artery and vein are ligated.
    *   A pyloromyotomy or pyloroplasty is often performed to facilitate gastric emptying, as vagal denervation will occur.
    *   Extensive lymphadenectomy of the celiac axis, lesser curvature, and greater curvature nodes is performed.
    *   A feeding jejunostomy tube is typically placed for postoperative enteral nutrition.

  \item  **Thoracic Phase (Right Thoracotomy or Thoracoscopy):**
    *   The patient is repositioned to the left lateral decubitus position.
    *   A right posterolateral thoracotomy or thoracoscopic ports are placed.
    *   The esophagus is dissected from its mediastinal attachments, from the diaphragm up to the thoracic inlet.
    *   Mediastinal lymphadenectomy, including subcarinal, paraesophageal, and paratracheal nodes, is performed.
    *   The esophagus is transected distally above the tumor and proximally at a level determined by tumor extent and oncologic margins.

  \item  **Cervical Phase (Left Cervical Incision):**
    *   The patient is repositioned supine with the neck extended and turned to the right.
    *   A left cervical incision is made along the anterior border of the sternocleidomastoid muscle.
    *   The cervical esophagus is identified and mobilized.
    *   The gastric conduit, prepared during the abdominal phase, is carefully pulled up through the posterior mediastinum (esophageal bed) into the neck.
    *   A tension-free esophagogastric anastomosis is created in the neck, typically end-to-side or end-to-end, using hand-sewn or stapled techniques.

  \item  **Closure and Drains:**
    *   Chest tubes are placed in the thoracic cavity to drain fluid and air.
    *   The abdominal, thoracic, and cervical incisions are closed in layers.
    *   A nasogastric tube is usually left in place for gastric decompression.
\end{enumerate}
Minimally invasive esophagectomy (MIE) utilizes laparoscopic and thoracoscopic techniques to perform these steps, aiming to reduce surgical trauma, pain, and pulmonary complications, while maintaining the oncologic principles of open surgery. Robotic-assisted MIE further enhances precision and visualization.

\section*{Preoperative Workup \& Preparation}
Thorough preoperative workup and preparation are paramount for patients undergoing esophagectomy, given the magnitude of the surgery and the often-compromised health status of patients with esophageal cancer. This includes comprehensive medical and cardiac evaluation to assess fitness for surgery, often involving an electrocardiogram, echocardiogram, and pulmonary function tests. Laboratory studies include a complete blood count, comprehensive metabolic panel, coagulation profile, and nutritional markers such as albumin and prealbumin. Imaging for staging is critical and typically involves a computed tomography (CT) scan of the chest, abdomen, and pelvis, endoscopic ultrasound (EUS) with fine-needle aspiration for local staging and nodal assessment, and often a positron emission tomography (PET) scan to detect distant metastases.

\begin{itemize}
  \item **Nutritional Optimization:** Many patients are malnourished; aggressive nutritional support, sometimes via a preoperative feeding jejunostomy or total parenteral nutrition, is crucial.
  \item **Cardiopulmonary Optimization:** Smoking cessation is mandatory. Pulmonary rehabilitation, bronchodilators for chronic obstructive pulmonary disease, and management of cardiac comorbidities are essential.
  \item **Dental Evaluation:** To minimize the risk of aspiration pneumonia from poor oral hygiene.
  \item **Anesthesia Clearance:** A thorough assessment by an anesthesiologist to determine suitability for prolonged general anesthesia and single-lung ventilation.
  \item **Medication Adjustment:** Anticoagulants are typically held, and other medications adjusted as per protocol.
  \item **Nil Per Os (NPO):** Patients are kept NPO for at least 6-8 hours prior to surgery.
  \item **Prophylactic Antibiotics:** Administered intravenously within one hour of incision.
  \item **Informed Consent:** Detailed discussion with the patient and family regarding the procedure, potential risks, expected recovery, and long-term implications.
\end{itemize}

\section*{Contraindications \& Precautions}
Esophagectomy is a high-risk procedure, and careful patient selection is critical.

**Absolute Contraindications:**
\begin{itemize}
  \item Documented distant metastatic disease (e.g., liver, lung, bone, non-regional lymph nodes) that is not amenable to curative resection.
  \item Extensive local invasion of vital structures such as the aorta, trachea, or major airways, making R0 resection impossible without unacceptable morbidity or mortality.
  \item Severe, uncorrectable medical comorbidities (e.g., severe cardiac dysfunction, end-stage lung disease, uncontrolled liver cirrhosis) that render the patient unfit to tolerate prolonged general anesthesia and major surgery.
  \item Patient refusal or inability to provide informed consent.
\end{itemize}

**Relative Contraindications \& Precautions:**
\begin{itemize}
  \item Advanced age (typically >80 years) with poor functional status or significant comorbidities, where the risks may outweigh the benefits.
  \item Significant malnutrition or cachexia, which increases the risk of complications; aggressive preoperative nutritional support is a precaution.
  \item Severe cardiopulmonary disease that is borderline for surgical tolerance; meticulous optimization and careful intraoperative management are crucial precautions.
  \item Prior extensive abdominal or thoracic surgery that may complicate dissection and conduit mobilization.
  \item Active infection or sepsis, which should be treated and resolved before elective surgery.
  \item Renal insufficiency or diabetes mellitus, requiring careful perioperative management to prevent complications.
\end{itemize}
Precautions also include meticulous surgical technique to minimize blood loss, prevent nerve injury (especially recurrent laryngeal), and ensure a well-perfused, tension-free anastomosis. Postoperative vigilance for anastomotic leak and pulmonary complications is paramount.

\section*{Complications \& Risks}
Esophagectomy is associated with a significant risk of complications, both general surgical and procedure-specific.

**General Surgical Complications:**
\begin{itemize}
  \item **Bleeding and Hemorrhage:** Intraoperative or postoperative, potentially requiring transfusion or reoperation.
  \item **Infection:** Surgical site infection, pneumonia, urinary tract infection, or sepsis.
  \item **Anesthesia-related risks:** Adverse reactions to medications, cardiac events, stroke, nerve damage.
  \item **Deep Vein Thrombosis (DVT) and Pulmonary Embolism (PE):** Due to prolonged immobility.
  \item **Cardiac Complications:** Atrial fibrillation, myocardial infarction, heart failure.
\end{itemize}

**Procedure-Specific Complications:**
\begin{itemize}
  \item **Anastomotic Leakage:** The most feared complication, occurring in 10-20\% of cases, leading to mediastinitis, empyema, sepsis, and high mortality.
  \item **Pulmonary Complications:** Pneumonia, atelectasis, acute respiratory distress syndrome (ARDS), pleural effusion, and respiratory failure are common due especially to single-lung ventilation and proximity to the lungs.
  \item **Recurrent Laryngeal Nerve (RLN) Injury:** Can cause vocal cord paralysis, hoarseness, and dysphagia, increasing the risk of aspiration.
  \item **Chylothorax:** Injury to the thoracic duct, leading to leakage of lymphatic fluid into the pleural space, requiring drainage and potentially reoperation or dietary modification.
  \item **Gastric Conduit Ischemia/Necrosis:** Inadequate blood supply to the mobilized stomach, leading to conduit breakdown and catastrophic complications.
  \item **Delayed Gastric Emptying:** Due to vagal nerve transection, often requiring prokinetic agents or long-term dietary modifications.
  \item **Anastomotic Stricture:** Narrowing at the anastomotic site, causing dysphagia, often requiring endoscopic dilation.
  \item **Reflux Symptoms:** Despite reconstruction, patients may experience bile reflux or symptoms similar to GERD.
  \item **Post-esophagectomy Syndrome:** A constellation of symptoms including dumping syndrome, early satiety, and weight loss.
  \item **Perioperative Mortality:** Historically high, but now typically 2-5\% in high-volume centers.
\end{itemize}

\section*{Postoperative Recovery Timeline}
Immediately following esophagectomy, patients are transferred to the intensive care unit (ICU) for close monitoring. In the first 24-48 hours, focus is on hemodynamic stability, respiratory support (often mechanical ventilation initially), pain control, and vigilant observation for signs of anastomotic leak or other acute complications. Chest tubes are monitored for output, and the nasogastric tube is kept on low suction for gastric decompression. Early mobilization, often within 24 hours, is encouraged to prevent pulmonary complications.

Over the first 1-2 weeks, patients are gradually weaned from mechanical ventilation, chest tubes are removed once drainage subsides and no air leak is present, and the nasogastric tube is typically removed after a few days. Nutrition is initially provided via the feeding jejunostomy tube. Oral intake usually begins with sips of water or clear liquids after a contrast swallow study confirms anastomotic integrity, typically around postoperative day 5-7. If the swallow study is normal, the diet is slowly advanced to pureed foods and then soft solids. Patients are encouraged to ambulate frequently and participate in pulmonary rehabilitation. The average hospital stay ranges from 10 to 14 days, depending on the patient's recovery trajectory and absence of complications. Long-term recovery involves adapting to a new eating pattern (smaller, more frequent meals), managing potential reflux or dumping symptoms, and gradually increasing activity levels over several months. Full recovery of strength and stamina can take 3-6 months.

\section*{Surgical Findings \& Pathology}
During esophagectomy, the surgeon meticulously documents the location and size of the primary tumor, its macroscopic invasion into surrounding structures, the extent of lymphadenopathy, and the presence of any suspicious lesions in the resected specimen or surrounding tissues. The completeness of resection, particularly the proximal and distal margins, is visually assessed. The resected esophagus, stomach (if used as conduit), and all removed lymph nodes are sent for detailed pathological examination. The pathology report provides crucial information: the exact tumor type (e.g., adenocarcinoma, squamous cell carcinoma), histological grade, depth of invasion (T-stage), number and location of involved lymph nodes (N-stage), and the status of the resection margins (R0 for negative, R1 for microscopic positive, R2 for macroscopic positive). It also details the response to neoadjuvant therapy, if applicable, by assessing the percentage of residual tumor. These findings are critical for definitive staging (pathological stage), guiding adjuvant therapy decisions, and determining prognosis.

\section{Expected Postoperative Course}
A normal and successful postoperative course following esophagectomy is characterized by a gradual and uneventful recovery without major complications. Patients will experience incisional pain, which is managed with multimodal analgesia, progressively decreasing over the first few weeks. The ability to swallow liquids and then soft foods returns, and patients learn to adapt to their new anatomy, typically eating smaller, more frequent meals to avoid discomfort and manage early satiety. The feeding jejunostomy provides nutritional support until oral intake is sufficient, and it is usually removed a few weeks after discharge. Patients regain their mobility and strength, gradually resuming daily activities. The surgical wounds heal without infection, and the chest tubes are removed without prolonged air leaks or effusions. Ultimately, a normal course means the patient is discharged home, able to maintain adequate nutrition orally, and is free from significant complications such as anastomotic leak, severe pulmonary issues, or recurrent laryngeal nerve injury, allowing for a smoother transition to long-term recovery and potential adjuvant therapies.

\section{Postoperative Care \& Follow-up}
Postoperative care after esophagectomy is extensive and continues for many years.

\begin{itemize}
  \item **Hospital Discharge:** Patients are discharged once they are medically stable, pain is controlled, oral intake is adequate, and they or their caregivers are comfortable managing the feeding jejunostomy (if still present) and wound care.
  \item **First Follow-up Visit (2-4 weeks post-discharge):** To assess wound healing, remove any remaining sutures/staples, review pathology results, and discuss adjuvant therapy plans.
  \item **Nutritional Support:** Ongoing dietary counseling is essential to manage the altered digestive physiology, including advice on small, frequent meals, avoiding certain foods, and managing reflux or dumping symptoms. Nutritional supplements may be prescribed.
  \item **Physical Therapy/Rehabilitation:** Encouraged to regain strength, improve lung function, and return to baseline activity levels.
  \item **Surveillance for Recurrence:** Regular follow-up with imaging (CT scans of chest/abdomen/pelvis) and endoscopy is crucial for early detection of recurrent disease, typically every 3-6 months for the first 2-3 years, then annually.
  \item **Management of Late Complications:** Surveillance for and endoscopic dilation of anastomotic strictures, management of chronic reflux, and addressing any persistent dysphagia or weight loss.
  \item **Warning Signs:** Patients are educated to contact their surgeon immediately for symptoms such as persistent fever, increasing pain, difficulty swallowing, persistent nausea/vomiting, significant weight loss, or signs of wound infection.
\item **Psychological Support:** Many patients benefit from support groups or counseling to cope with the physical and emotional challenges of cancer and major surgery.
\end{itemize}

\section*{Epidemiology}
Esophageal cancer is the eighth most common cancer worldwide and the sixth leading cause of cancer-related death, with significant geographical variation in incidence and histological type. Squamous cell carcinoma (SCC) is historically more prevalent globally, particularly in parts of Asia and Africa, and is strongly associated with tobacco smoking, alcohol consumption, hot beverages, and nutritional deficiencies. In Western countries, however, the incidence of esophageal adenocarcinoma (EAC) has risen dramatically over the past few decades, now surpassing SCC in many regions. EAC typically arises from Barrett's esophagus, a complication of chronic gastroesophageal reflux disease (GERD), and is strongly linked to obesity. Esophageal cancer predominantly affects men and older adults, with the median age at diagnosis typically in the sixth or seventh decade of life. Despite advances in multimodal therapy, the overall prognosis remains poor due to late presentation and aggressive tumor biology, with 5-year survival rates ranging from 15-40\% depending on stage. Perioperative mortality for esophagectomy, while significantly reduced from historical figures, still ranges from 2-5\% in high-volume centers.

\section*{Outcomes \& Prognosis}
The outcomes and prognosis following esophagectomy are primarily dictated by the pathological stage of the cancer, particularly the completeness of resection (R0 status) and the extent of lymph node involvement. Patients with early-stage disease (T1-T2, N0) and a complete R0 resection have the most favorable prognosis, with 5-year survival rates potentially exceeding 50-60\%. Conversely, the presence of positive lymph nodes (N1 or higher) or positive surgical margins (R1/R2) significantly worsens prognosis, increasing the risk of recurrence and reducing long-term survival. The landmark CROSS trial demonstrated that neoadjuvant chemoradiation followed by surgery significantly improved overall survival compared to surgery alone for locally advanced esophageal cancer. Functional outcomes are generally good, with most patients able to resume oral intake, although long-term quality of life can be affected by symptoms such as dysphagia, reflux, dumping syndrome, and weight loss. While perioperative mortality has decreased, morbidity remains high, with pulmonary complications and anastomotic leaks being major contributors. Recurrence rates vary but can be as high as 30-50\% even after R0 resection, necessitating vigilant long-term surveillance. Prognostic factors also include tumor biology, patient comorbidities, and the volume and expertise of the treating center.

\section*{References}
\begin{enumerate}
  \item MedlinePlus. Esophagectomy. U.S. National Library of Medicine; 2024. Available from: https://medlineplus.gov/ency/article/002947.htm
  \item Braghetto I, Csendes A, Burdiles P, et al. Transhiatal esophagectomy: a 25-year experience. World J Surg. 2005;29(11):1405-1410.
  \item Schwartz's Principles of Surgery. 11th ed. New York, NY: McGraw-Hill Education; 2019. Chapter 29: Esophagus.
  \item van Hagen P, Hulshof MC, van Lanschot JJ, et al. Preoperative chemoradiotherapy for esophageal or junctional cancer. N Engl J Med. 2012;366(22):2074-2084.
\end{enumerate}

\clearpage